\section{Discussion}

The results link firm size to three related outcomes: wage inequality,
labor-market segmentation, and possible allocative frictions. The inequality
result is direct: workers in larger firm-size categories earn more than workers
in solo work and micro-firms, both in raw averages and after conditioning on
observed characteristics. The segmentation result comes from the role of
formality. Adding formality to the main regression absorbs part of the raw
large-firm premium, and formal workers earn more than informal workers within
every firm-size category. At the same time, the premium does not disappear after
controlling for formality, and the size profile differs sharply between formal
and informal employment.

The formality results narrow the policy-relevant interpretation. Formal employment raises wage levels within every firm-size category, but the conditional size premium is steepest among informal workers. Within formal employment, the premium relative to formal solo workers is small and not statistically significant, although workers in 101+ firms earn more than workers in intermediate-size formal firms. If the conditional wage gaps partly reflect productivity-related frictions rather than only unobserved sorting, the evidence points to three margins for possible productivity gains: moving workers and jobs from informality to formality, scaling informal micro-units into larger and more organized firms, and, more selectively, growth or movement within the formal sector from intermediate-size firms toward larger employers.

These interpretations must remain cautious, though. A wage premium that survives observable controls does not identify a single mechanism as an explanation. As \citet{BrownMedoff1989} and \citet{OiIdson1999} emphasize, the remaining gap may reflect unobserved worker quality, firm productivity, rent sharing, efficiency wages, or match quality. The developing-country evidence cautions against reducing the premium to sorting alone: \citet{SoderbomTealWambugu2005} find that the firm-size effect remains substantial in Kenya and Ghana after controlling for individual fixed effects. Matched employer-employee studies also show that employer-associated characteristics matter for wage variation \citep{AbowdKramarzMargolis1999,CardHeiningKline2013,Song2019}, and related evidence from Brazil reaches a similar conclusion for Latin America \citep{AlvarezBenguriaEngbomMoser2018,GerardLagosSeverniniCard2021, EngbomMoser2022}. Recent German evidence sharpens this caution: \citet{DiegmannMullerSchoefer2026} find that the employer-size wage effect disappears when establishment size is compared within the same multi-establishment firm, suggesting that the premium may reflect employer-wide wage policies, rents, or other firm attributes rather than establishment size per se. The elasticity benchmark in Section~\ref{subsec:elasticity-benchmark} is consistent with a large Colombian size-wage association, especially among informal workers, but GEIH does not identify employers or follow workers across firms. The paper therefore cannot separate worker sorting from other competing mechanisms, so the estimated premiums and elasticities should be interpreted as diagnostic magnitudes rather than causal effects of firm size.

The evidence is also not a direct test of misallocation. Misallocation in the strict sense requires comparing marginal products across firms or sectors, as in \citet{hsieh2009misallocation}. The link with allocative efficiency comes from the marginal-product condition: in an efficient allocation, labor should move across firms until the marginal product of labor is equalized across uses. In the competitive benchmark, wages reflect the value of marginal product, so persistent wage differences across firm sizes are informative about possible allocative frictions. If larger firms pay more than smaller firms for workers with the same observable characteristics, this may indicate that labor is more valuable in larger firms or that barriers prevent workers from moving toward higher-productivity matches. The evidence is only a marker, not proof, because wages may also reflect unobserved worker quality or other unobserved features of the employment relationship. Testing misallocation directly would require firm-level data on output, capital, and employment, which would allow us to estimate marginal products rather than infer them indirectly from wages.

The policy implication is therefore not that policy should favor large firms as such or push workers mechanically into large firms. The relevant objective is to reduce barriers that keep productive firms small, informal, or disconnected from larger markets, while also improving workers' access to better firm matches. The formality results sharpen this point. Because formal employment raises wage levels within every firm-size category, moving jobs from informality to formality remains a policy priority. Because the firm-size premium is steepest among informal workers, a second policy priority is helping productive informal micro-units become larger, more organized, and more connected to markets. This support should be designed as a path toward formalization, not as a subsidy for remaining informal; otherwise, it could weaken incentives for formal firms to remain formal. The evidence gives less support to a policy focus on scaling up formal firms: within formal employment, the premium relative to formal solo workers is small, although workers in 101+ firms earn more than workers in intermediate-size formal firms. This pattern is consistent with a segmented labor market in which the strongest wage-based signs of allocative frictions are concentrated in informal employment. The largest efficiency gains therefore appear to lie in formalization and in the growth of productive informal firms, rather than in firm growth within the formal sector.
